%=============================================================
%  L01 ARTIFACT SUITE — MASTER MANIFEST
%  Module I.3, Lecture L01: Pure/Mixed States and the Density Matrix
%  QM Programme
%=============================================================
\documentclass[11pt,a4paper]{article}
%=============================================================
%  QM PROGRAMME — SHARED PREAMBLE
%  Paste %=============================================================
%  QM PROGRAMME — SHARED PREAMBLE
%  Paste %=============================================================
%  QM PROGRAMME — SHARED PREAMBLE
%  Paste \input{preamble} at the top of each artifact file,
%  or compile standalone by wrapping in \documentclass{article}
%=============================================================
\usepackage[T1]{fontenc}
\usepackage[utf8]{inputenc}
\usepackage[margin=2.5cm]{geometry}
\usepackage{amsmath,amssymb,amsthm}
\usepackage{physics}
\usepackage{braket}
\usepackage{booktabs,array,longtable,tabularx}
\usepackage{enumitem}
\usepackage{fancyhdr}
\usepackage[dvipsnames,svgnames,table]{xcolor}
\usepackage{hyperref}
\usepackage{parskip}
\usepackage{tcolorbox}
\tcbuselibrary{skins,breakable}
%\usepackage{microtype}  % disabled: font expansion incompatibility

%--- Programme colours ----------------------------------------
\definecolor{darkblue}{RGB}{10,40,90}
\definecolor{midblue}{RGB}{30,80,160}
\definecolor{lightblue}{RGB}{220,233,255}
\definecolor{accentgold}{RGB}{180,140,20}
\definecolor{lightgold}{RGB}{255,248,220}
\definecolor{darkgray}{RGB}{60,60,60}
\definecolor{tier1col}{RGB}{180,230,180}
\definecolor{tier2col}{RGB}{150,210,240}
\definecolor{tier3col}{RGB}{200,180,240}
\definecolor{tier4col}{RGB}{255,210,150}
\definecolor{tier5col}{RGB}{255,170,150}
\definecolor{ugblue}{RGB}{25,90,170}
\definecolor{lightugblue}{RGB}{215,230,255}
\definecolor{uggreen}{RGB}{20,120,70}
\definecolor{lightuggreen}{RGB}{210,245,225}
\definecolor{gradred}{RGB}{160,30,30}
\definecolor{lightgradred}{RGB}{255,220,215}
\definecolor{gradorange}{RGB}{170,80,10}
\definecolor{lightgradorange}{RGB}{255,235,210}
\definecolor{respath}{RGB}{90,20,130}
\definecolor{lightrespath}{RGB}{240,220,255}
\definecolor{artpurple}{RGB}{100,30,130}
\definecolor{lightartpurple}{RGB}{235,215,245}

%--- tcolorbox styles -----------------------------------------
\tcbset{
  ugconceptbox/.style={enhanced,breakable,
    colback=lightugblue,colframe=ugblue,
    fonttitle=\bfseries\color{white},coltitle=white,
    attach boxed title to top left={yshift=-2mm,xshift=4mm},
    boxed title style={colback=ugblue,sharp corners,boxrule=0.5pt},
    sharp corners=south,boxrule=0.8pt,
    left=6pt,right=6pt,top=8pt,bottom=6pt},
  ugprobbox/.style={enhanced,breakable,
    colback=lightuggreen,colframe=uggreen,
    fonttitle=\bfseries\color{white},coltitle=white,
    attach boxed title to top left={yshift=-2mm,xshift=4mm},
    boxed title style={colback=uggreen,sharp corners,boxrule=0.5pt},
    sharp corners=south,boxrule=0.8pt,
    left=6pt,right=6pt,top=8pt,bottom=6pt},
  gradconceptbox/.style={enhanced,breakable,
    colback=lightgradred,colframe=gradred,
    fonttitle=\bfseries\color{white},coltitle=white,
    attach boxed title to top left={yshift=-2mm,xshift=4mm},
    boxed title style={colback=gradred,sharp corners,boxrule=0.5pt},
    sharp corners=south,boxrule=0.8pt,
    left=6pt,right=6pt,top=8pt,bottom=6pt},
  gradprobbox/.style={enhanced,breakable,
    colback=lightgradorange,colframe=gradorange,
    fonttitle=\bfseries\color{white},coltitle=white,
    attach boxed title to top left={yshift=-2mm,xshift=4mm},
    boxed title style={colback=gradorange,sharp corners,boxrule=0.5pt},
    sharp corners=south,boxrule=0.8pt,
    left=6pt,right=6pt,top=8pt,bottom=6pt},
  researchbox/.style={enhanced,breakable,
    colback=lightrespath,colframe=respath,
    fonttitle=\bfseries\color{white},coltitle=white,
    attach boxed title to top left={yshift=-2mm,xshift=4mm},
    boxed title style={colback=respath,sharp corners,boxrule=0.5pt},
    sharp corners=south,boxrule=0.8pt,
    left=6pt,right=6pt,top=8pt,bottom=6pt},
  rubricbox/.style={enhanced,
    colback=lightartpurple,colframe=artpurple,
    fonttitle=\bfseries\color{white},coltitle=white,
    attach boxed title to top left={yshift=-2mm,xshift=4mm},
    boxed title style={colback=artpurple,sharp corners,boxrule=0.5pt},
    sharp corners=south,boxrule=0.7pt,
    left=5pt,right=5pt,top=6pt,bottom=5pt},
  goldbox/.style={enhanced,breakable,
    colback=lightgold,colframe=accentgold,
    boxrule=0.8pt,left=6pt,right=6pt,top=6pt,bottom=6pt,
    fontupper=\small},
  notebox/.style={enhanced,
    colback=lightblue,colframe=midblue,
    boxrule=0.6pt,rounded corners,
    left=5pt,right=5pt,top=4pt,bottom=4pt,
    fontupper=\small\itshape},
  trackheaderUG/.style={enhanced,
    colback=ugblue!12!white,colframe=ugblue,
    boxrule=1.2pt,sharp corners,
    left=6pt,right=6pt,top=5pt,bottom=5pt,
    fontupper=\bfseries\large\color{ugblue}},
  trackheaderGRAD/.style={enhanced,
    colback=gradred!12!white,colframe=gradred,
    boxrule=1.2pt,sharp corners,
    left=6pt,right=6pt,top=5pt,bottom=5pt,
    fontupper=\bfseries\large\color{gradred}},
  trackheaderRES/.style={enhanced,
    colback=respath!12!white,colframe=respath,
    boxrule=1.2pt,sharp corners,
    left=6pt,right=6pt,top=5pt,bottom=5pt,
    fontupper=\bfseries\large\color{respath}},
}

%--- Programme macros -----------------------------------------
\newcommand{\instructornote}[1]{%
  \begin{tcolorbox}[notebox]\textbf{Instructor note:} #1\end{tcolorbox}}
\newcommand{\tiertag}[2]{%
  \colorbox{#1}{\small\bfseries\color{darkblue}\ #2\ }}
\newcommand{\essential}{$\bigstar\bigstar\bigstar$}
\newcommand{\recommended}{$\bigstar\bigstar$}
\newcommand{\useful}{$\bigstar$}
\newcommand{\rhat}{\hat{\rho}}
\newcommand{\Ahat}{\hat{A}}
\newcommand{\Hhat}{\hat{H}}
\newcommand{\Uhat}{\hat{U}}
\newcommand{\Ihat}{\hat{\mathbb{1}}}
\newcommand{\TrB}{\operatorname{Tr}_B}
\newcommand{\psiket}{\ket{\psi}}
\newcommand{\phiket}{\ket{\phi}}
 at the top of each artifact file,
%  or compile standalone by wrapping in \documentclass{article}
%=============================================================
\usepackage[T1]{fontenc}
\usepackage[utf8]{inputenc}
\usepackage[margin=2.5cm]{geometry}
\usepackage{amsmath,amssymb,amsthm}
\usepackage{physics}
\usepackage{braket}
\usepackage{booktabs,array,longtable,tabularx}
\usepackage{enumitem}
\usepackage{fancyhdr}
\usepackage[dvipsnames,svgnames,table]{xcolor}
\usepackage{hyperref}
\usepackage{parskip}
\usepackage{tcolorbox}
\tcbuselibrary{skins,breakable}
%\usepackage{microtype}  % disabled: font expansion incompatibility

%--- Programme colours ----------------------------------------
\definecolor{darkblue}{RGB}{10,40,90}
\definecolor{midblue}{RGB}{30,80,160}
\definecolor{lightblue}{RGB}{220,233,255}
\definecolor{accentgold}{RGB}{180,140,20}
\definecolor{lightgold}{RGB}{255,248,220}
\definecolor{darkgray}{RGB}{60,60,60}
\definecolor{tier1col}{RGB}{180,230,180}
\definecolor{tier2col}{RGB}{150,210,240}
\definecolor{tier3col}{RGB}{200,180,240}
\definecolor{tier4col}{RGB}{255,210,150}
\definecolor{tier5col}{RGB}{255,170,150}
\definecolor{ugblue}{RGB}{25,90,170}
\definecolor{lightugblue}{RGB}{215,230,255}
\definecolor{uggreen}{RGB}{20,120,70}
\definecolor{lightuggreen}{RGB}{210,245,225}
\definecolor{gradred}{RGB}{160,30,30}
\definecolor{lightgradred}{RGB}{255,220,215}
\definecolor{gradorange}{RGB}{170,80,10}
\definecolor{lightgradorange}{RGB}{255,235,210}
\definecolor{respath}{RGB}{90,20,130}
\definecolor{lightrespath}{RGB}{240,220,255}
\definecolor{artpurple}{RGB}{100,30,130}
\definecolor{lightartpurple}{RGB}{235,215,245}

%--- tcolorbox styles -----------------------------------------
\tcbset{
  ugconceptbox/.style={enhanced,breakable,
    colback=lightugblue,colframe=ugblue,
    fonttitle=\bfseries\color{white},coltitle=white,
    attach boxed title to top left={yshift=-2mm,xshift=4mm},
    boxed title style={colback=ugblue,sharp corners,boxrule=0.5pt},
    sharp corners=south,boxrule=0.8pt,
    left=6pt,right=6pt,top=8pt,bottom=6pt},
  ugprobbox/.style={enhanced,breakable,
    colback=lightuggreen,colframe=uggreen,
    fonttitle=\bfseries\color{white},coltitle=white,
    attach boxed title to top left={yshift=-2mm,xshift=4mm},
    boxed title style={colback=uggreen,sharp corners,boxrule=0.5pt},
    sharp corners=south,boxrule=0.8pt,
    left=6pt,right=6pt,top=8pt,bottom=6pt},
  gradconceptbox/.style={enhanced,breakable,
    colback=lightgradred,colframe=gradred,
    fonttitle=\bfseries\color{white},coltitle=white,
    attach boxed title to top left={yshift=-2mm,xshift=4mm},
    boxed title style={colback=gradred,sharp corners,boxrule=0.5pt},
    sharp corners=south,boxrule=0.8pt,
    left=6pt,right=6pt,top=8pt,bottom=6pt},
  gradprobbox/.style={enhanced,breakable,
    colback=lightgradorange,colframe=gradorange,
    fonttitle=\bfseries\color{white},coltitle=white,
    attach boxed title to top left={yshift=-2mm,xshift=4mm},
    boxed title style={colback=gradorange,sharp corners,boxrule=0.5pt},
    sharp corners=south,boxrule=0.8pt,
    left=6pt,right=6pt,top=8pt,bottom=6pt},
  researchbox/.style={enhanced,breakable,
    colback=lightrespath,colframe=respath,
    fonttitle=\bfseries\color{white},coltitle=white,
    attach boxed title to top left={yshift=-2mm,xshift=4mm},
    boxed title style={colback=respath,sharp corners,boxrule=0.5pt},
    sharp corners=south,boxrule=0.8pt,
    left=6pt,right=6pt,top=8pt,bottom=6pt},
  rubricbox/.style={enhanced,
    colback=lightartpurple,colframe=artpurple,
    fonttitle=\bfseries\color{white},coltitle=white,
    attach boxed title to top left={yshift=-2mm,xshift=4mm},
    boxed title style={colback=artpurple,sharp corners,boxrule=0.5pt},
    sharp corners=south,boxrule=0.7pt,
    left=5pt,right=5pt,top=6pt,bottom=5pt},
  goldbox/.style={enhanced,breakable,
    colback=lightgold,colframe=accentgold,
    boxrule=0.8pt,left=6pt,right=6pt,top=6pt,bottom=6pt,
    fontupper=\small},
  notebox/.style={enhanced,
    colback=lightblue,colframe=midblue,
    boxrule=0.6pt,rounded corners,
    left=5pt,right=5pt,top=4pt,bottom=4pt,
    fontupper=\small\itshape},
  trackheaderUG/.style={enhanced,
    colback=ugblue!12!white,colframe=ugblue,
    boxrule=1.2pt,sharp corners,
    left=6pt,right=6pt,top=5pt,bottom=5pt,
    fontupper=\bfseries\large\color{ugblue}},
  trackheaderGRAD/.style={enhanced,
    colback=gradred!12!white,colframe=gradred,
    boxrule=1.2pt,sharp corners,
    left=6pt,right=6pt,top=5pt,bottom=5pt,
    fontupper=\bfseries\large\color{gradred}},
  trackheaderRES/.style={enhanced,
    colback=respath!12!white,colframe=respath,
    boxrule=1.2pt,sharp corners,
    left=6pt,right=6pt,top=5pt,bottom=5pt,
    fontupper=\bfseries\large\color{respath}},
}

%--- Programme macros -----------------------------------------
\newcommand{\instructornote}[1]{%
  \begin{tcolorbox}[notebox]\textbf{Instructor note:} #1\end{tcolorbox}}
\newcommand{\tiertag}[2]{%
  \colorbox{#1}{\small\bfseries\color{darkblue}\ #2\ }}
\newcommand{\essential}{$\bigstar\bigstar\bigstar$}
\newcommand{\recommended}{$\bigstar\bigstar$}
\newcommand{\useful}{$\bigstar$}
\newcommand{\rhat}{\hat{\rho}}
\newcommand{\Ahat}{\hat{A}}
\newcommand{\Hhat}{\hat{H}}
\newcommand{\Uhat}{\hat{U}}
\newcommand{\Ihat}{\hat{\mathbb{1}}}
\newcommand{\TrB}{\operatorname{Tr}_B}
\newcommand{\psiket}{\ket{\psi}}
\newcommand{\phiket}{\ket{\phi}}
 at the top of each artifact file,
%  or compile standalone by wrapping in \documentclass{article}
%=============================================================
\usepackage[T1]{fontenc}
\usepackage[utf8]{inputenc}
\usepackage[margin=2.5cm]{geometry}
\usepackage{amsmath,amssymb,amsthm}
\usepackage{physics}
\usepackage{braket}
\usepackage{booktabs,array,longtable,tabularx}
\usepackage{enumitem}
\usepackage{fancyhdr}
\usepackage[dvipsnames,svgnames,table]{xcolor}
\usepackage{hyperref}
\usepackage{parskip}
\usepackage{tcolorbox}
\tcbuselibrary{skins,breakable}
%\usepackage{microtype}  % disabled: font expansion incompatibility

%--- Programme colours ----------------------------------------
\definecolor{darkblue}{RGB}{10,40,90}
\definecolor{midblue}{RGB}{30,80,160}
\definecolor{lightblue}{RGB}{220,233,255}
\definecolor{accentgold}{RGB}{180,140,20}
\definecolor{lightgold}{RGB}{255,248,220}
\definecolor{darkgray}{RGB}{60,60,60}
\definecolor{tier1col}{RGB}{180,230,180}
\definecolor{tier2col}{RGB}{150,210,240}
\definecolor{tier3col}{RGB}{200,180,240}
\definecolor{tier4col}{RGB}{255,210,150}
\definecolor{tier5col}{RGB}{255,170,150}
\definecolor{ugblue}{RGB}{25,90,170}
\definecolor{lightugblue}{RGB}{215,230,255}
\definecolor{uggreen}{RGB}{20,120,70}
\definecolor{lightuggreen}{RGB}{210,245,225}
\definecolor{gradred}{RGB}{160,30,30}
\definecolor{lightgradred}{RGB}{255,220,215}
\definecolor{gradorange}{RGB}{170,80,10}
\definecolor{lightgradorange}{RGB}{255,235,210}
\definecolor{respath}{RGB}{90,20,130}
\definecolor{lightrespath}{RGB}{240,220,255}
\definecolor{artpurple}{RGB}{100,30,130}
\definecolor{lightartpurple}{RGB}{235,215,245}

%--- tcolorbox styles -----------------------------------------
\tcbset{
  ugconceptbox/.style={enhanced,breakable,
    colback=lightugblue,colframe=ugblue,
    fonttitle=\bfseries\color{white},coltitle=white,
    attach boxed title to top left={yshift=-2mm,xshift=4mm},
    boxed title style={colback=ugblue,sharp corners,boxrule=0.5pt},
    sharp corners=south,boxrule=0.8pt,
    left=6pt,right=6pt,top=8pt,bottom=6pt},
  ugprobbox/.style={enhanced,breakable,
    colback=lightuggreen,colframe=uggreen,
    fonttitle=\bfseries\color{white},coltitle=white,
    attach boxed title to top left={yshift=-2mm,xshift=4mm},
    boxed title style={colback=uggreen,sharp corners,boxrule=0.5pt},
    sharp corners=south,boxrule=0.8pt,
    left=6pt,right=6pt,top=8pt,bottom=6pt},
  gradconceptbox/.style={enhanced,breakable,
    colback=lightgradred,colframe=gradred,
    fonttitle=\bfseries\color{white},coltitle=white,
    attach boxed title to top left={yshift=-2mm,xshift=4mm},
    boxed title style={colback=gradred,sharp corners,boxrule=0.5pt},
    sharp corners=south,boxrule=0.8pt,
    left=6pt,right=6pt,top=8pt,bottom=6pt},
  gradprobbox/.style={enhanced,breakable,
    colback=lightgradorange,colframe=gradorange,
    fonttitle=\bfseries\color{white},coltitle=white,
    attach boxed title to top left={yshift=-2mm,xshift=4mm},
    boxed title style={colback=gradorange,sharp corners,boxrule=0.5pt},
    sharp corners=south,boxrule=0.8pt,
    left=6pt,right=6pt,top=8pt,bottom=6pt},
  researchbox/.style={enhanced,breakable,
    colback=lightrespath,colframe=respath,
    fonttitle=\bfseries\color{white},coltitle=white,
    attach boxed title to top left={yshift=-2mm,xshift=4mm},
    boxed title style={colback=respath,sharp corners,boxrule=0.5pt},
    sharp corners=south,boxrule=0.8pt,
    left=6pt,right=6pt,top=8pt,bottom=6pt},
  rubricbox/.style={enhanced,
    colback=lightartpurple,colframe=artpurple,
    fonttitle=\bfseries\color{white},coltitle=white,
    attach boxed title to top left={yshift=-2mm,xshift=4mm},
    boxed title style={colback=artpurple,sharp corners,boxrule=0.5pt},
    sharp corners=south,boxrule=0.7pt,
    left=5pt,right=5pt,top=6pt,bottom=5pt},
  goldbox/.style={enhanced,breakable,
    colback=lightgold,colframe=accentgold,
    boxrule=0.8pt,left=6pt,right=6pt,top=6pt,bottom=6pt,
    fontupper=\small},
  notebox/.style={enhanced,
    colback=lightblue,colframe=midblue,
    boxrule=0.6pt,rounded corners,
    left=5pt,right=5pt,top=4pt,bottom=4pt,
    fontupper=\small\itshape},
  trackheaderUG/.style={enhanced,
    colback=ugblue!12!white,colframe=ugblue,
    boxrule=1.2pt,sharp corners,
    left=6pt,right=6pt,top=5pt,bottom=5pt,
    fontupper=\bfseries\large\color{ugblue}},
  trackheaderGRAD/.style={enhanced,
    colback=gradred!12!white,colframe=gradred,
    boxrule=1.2pt,sharp corners,
    left=6pt,right=6pt,top=5pt,bottom=5pt,
    fontupper=\bfseries\large\color{gradred}},
  trackheaderRES/.style={enhanced,
    colback=respath!12!white,colframe=respath,
    boxrule=1.2pt,sharp corners,
    left=6pt,right=6pt,top=5pt,bottom=5pt,
    fontupper=\bfseries\large\color{respath}},
}

%--- Programme macros -----------------------------------------
\newcommand{\instructornote}[1]{%
  \begin{tcolorbox}[notebox]\textbf{Instructor note:} #1\end{tcolorbox}}
\newcommand{\tiertag}[2]{%
  \colorbox{#1}{\small\bfseries\color{darkblue}\ #2\ }}
\newcommand{\essential}{$\bigstar\bigstar\bigstar$}
\newcommand{\recommended}{$\bigstar\bigstar$}
\newcommand{\useful}{$\bigstar$}
\newcommand{\rhat}{\hat{\rho}}
\newcommand{\Ahat}{\hat{A}}
\newcommand{\Hhat}{\hat{H}}
\newcommand{\Uhat}{\hat{U}}
\newcommand{\Ihat}{\hat{\mathbb{1}}}
\newcommand{\TrB}{\operatorname{Tr}_B}
\newcommand{\psiket}{\ket{\psi}}
\newcommand{\phiket}{\ket{\phi}}


\pagestyle{fancy}
\fancyhf{}
\fancyhead[L]{\small\color{darkgray}\textit{L01 --- Master Manifest}}
\fancyhead[R]{\small\color{darkgray}\thepage}
\fancyfoot[C]{\small\color{darkgray}\textit{L01 Artifact Suite --- Pure/Mixed States and the Density Matrix}}
\renewcommand{\headrulewidth}{0.4pt}
\hypersetup{colorlinks=true,linkcolor=midblue,pdftitle={L01-Manifest}}

\begin{document}

\begin{tcolorbox}[trackheaderUG]
L01 ARTIFACT SUITE --- Pure/Mixed States and the Density Matrix\\
\normalsize Module I.3 $\cdot$ QM Programme $\cdot$
Generated following \texttt{lecture\_generation\_prompt\_v3}
\end{tcolorbox}

\bigskip

%=============================================================
\section*{File Inventory}
%=============================================================

\subsection*{Lecture Body}

\begin{tabularx}{\textwidth}{llX}
\toprule
\textbf{File} & \textbf{Tiers} & \textbf{Contents}\\
\midrule
\texttt{lecture/L01\_lecture\_body.tex}
& All
& Metadata, pacing plan, full lecture exposition (Sections 3.1--3.4),
  connections (backward, forward, cross-module), and reading list
  with tier-tagged annotations.\\
\bottomrule
\end{tabularx}

\subsection*{Undergraduate Artifacts (UG Track)}

\begin{tabularx}{\textwidth}{llX}
\toprule
\textbf{File} & \textbf{Code} & \textbf{Contents / Scope}\\
\midrule
\texttt{ug/L01\_SCQU.tex}
& SCQU
& Simple Concept Questions UG. 10 questions: multiple choice (Q1--Q4),
  True/False + justify (Q5--Q8), short answer (Q9--Q10).
  Tiers 1--3. 15--20 min. 20 pts. Full rubric and model answers included.\\
\addlinespace
\texttt{ug/L01\_CCQU.tex}
& CCQU
& Complex Concept Questions UG. 8 questions:
  ``Explain why'' (Q1--Q4), ``What would change if'' (Q5--Q7),
  ``Identify the flaw'' (Q8--Q10). Tiers 2--3. 25--35 min. 20 pts.\\
\addlinespace
\texttt{ug/L01\_SPSU.tex}
& SPSU
& Simple Problem Set UG. 10 problems in three parts:
  Mechanics (A1--A4), Interpretation (B1--B3), Translation (C1--C3).
  Tiers 1--3. Rubric: Result 60\% / Method 30\% / Interpretation 10\%.\\
\addlinespace
\texttt{ug/L01\_CPSU.tex}
& CPSU
& Complex Problem Set UG. 5 multi-part extended proofs:
  (1) purity bounds and linear entropy; (2) Bloch-vector parameterisation;
  (3) partial trace and entanglement family; (4) ensemble ambiguity;
  (5) time evolution and purity conservation.
  Tier 3 primary. 3--6 hrs. 4-layer delivery system.\\
\addlinespace
\texttt{ug/L01\_SPJU.tex}
& SPJU
& Simple Projects UG. Three projects (choose one):
  (J1) Bloch sphere derivation and visualisation;
  (J2) numerical partial-trace explorer in Python;
  (J3) primary-source reading (Fano 1957) + NMR thermal state.
  Tiers 2--3. 1--2 weeks. 500--800 words or notebook.\\
\addlinespace
\texttt{ug/L01\_CPJU.tex}
& CPJU
& Complex Projects UG. Two projects (choose one):
  (CJ1) single-qubit quantum state tomography (theory + simulation);
  (CJ2) decoherence, pointer states, and the spin-boson model.
  Tier 3. 3--4 weeks. 2000--3000 words.\\
\bottomrule
\end{tabularx}

\subsection*{Graduate Artifacts (Grad Track)}

\begin{tabularx}{\textwidth}{llX}
\toprule
\textbf{File} & \textbf{Code} & \textbf{Contents / Scope}\\
\midrule
\texttt{grad/L01\_SCQG.tex}
& SCQG
& Simple Concept Questions Graduate. 10 questions:
  Structural probes (Q1--Q5: convexity, partial trace, relative entropy,
  sub-additivity) and functional-analytic probes (Q6--Q10: Uhlmann,
  Hilbert--Schmidt, trace-class, CP maps, Stinespring).
  Tiers 4--5. 20--30 min. 20 pts.\\
\addlinespace
\texttt{grad/L01\_CCQG.tex}
& CCQG
& Complex Concept Questions Graduate. 8 questions:
  ``Explain why'' (Q1--Q4: ensemble ambiguity, partial trace uniqueness,
  complete positivity, infinite-dimensional entropy);
  ``Compare and contrast'' (Q5--Q7: Shannon vs.\ von Neumann, pictures,
  Kraus vs.\ Stinespring); ``Identify and resolve'' (Q8: data-processing paradox).
  Tiers 4--5. 35--50 min.\\
\addlinespace
\texttt{grad/L01\_SPSG.tex}
& SPSG
& Simple Problem Set Graduate. 10 problems:
  Structural proofs (A1--A4: convexity, extremality, purity--rank, trace-class);
  Entropy and information (B1--B3: sub-additivity, concavity, purity--entropy);
  Channels and maps (C1--C3: depolarising, amplitude damping, positivity vs.\ CP).
  Tiers 4--5.\\
\addlinespace
\texttt{grad/L01\_CPSG.tex}
& CPSG
& Complex Problem Set Graduate. 5 extended proofs:
  (1) Uhlmann fidelity; (2) Schmidt decomposition and entanglement;
  (3) data-processing inequality; (4) GKSL derivation outline;
  (5) trace-class operators in $\ell^2(\mathbb{N})$.
  T4 Problems 1--3, T5 Problems 4--5.\\
\addlinespace
\texttt{grad/L01\_SPJG\_CPJG.tex}
& SPJG / CPJG
& Graduate Projects (simple and complex in one file).
  SPJG: three projects (SG1: SIC-POVM tomography; SG2: entanglement measures;
  SG3: Wigner function). 2--3 weeks.
  CPJG: five research-level projects (CG1: infinite-dim density operators;
  CG2: CPTP maps and GKSL; CG3: Petz recovery map; CG4: many-body RDMs;
  CG5: quantum de Finetti theorem). 4--6 weeks.
  Tiers 4--5.\\
\bottomrule
\end{tabularx}

\subsection*{Research Track Artifacts}

\begin{tabularx}{\textwidth}{llX}
\toprule
\textbf{File} & \textbf{Code} & \textbf{Contents / Scope}\\
\midrule
\texttt{research/L01\_WRQ\_ORQ.tex}
& WRQ / ORQ
& Written Research Questions (5, answer 2): (WRQ1) historical/conceptual:
  von Neumann's statistical operator; (WRQ2) foundational: density matrix
  and the measurement problem; (WRQ3) mathematical: Kraus representation;
  (WRQ4) information-theoretic: entropy and the second law;
  (WRQ5) applications: condensed matter and quantum chemistry.
  1000--1500 words each.\\
  & & Oral Research Questions (5, examiner selects 2): (ORQ1) density operator
  uniqueness; (ORQ2) partial trace vs.\ measurement; (ORQ3) complete positivity
  and entanglement; (ORQ4) GKSL and irreversibility; (ORQ5) entropy and holography.
  15-min viva format. Tier 5.\\
\bottomrule
\end{tabularx}

\subsection*{Bibliography}

\begin{tabularx}{\textwidth}{lX}
\toprule
\textbf{File} & \textbf{Contents}\\
\midrule
\texttt{bibliography/L01\_bibliography\_additions.tex}
& Full annotated bibliography for the L01 artifact suite.
  All references tier-tagged and cross-referenced to master bibliography
  (\texttt{qm\_bibliography.tex}). Six new entries identified for addition
  to the master bibliography (Zurek 2003; Breuer \& Petruccione 2002;
  Plenio \& Virmani 2007; Fawzi \& Renner 2015; Li \& Haldane 2008;
  White 1992).\\
\bottomrule
\end{tabularx}

\subsection*{Shared Infrastructure}

\begin{tabularx}{\textwidth}{lX}
\toprule
\textbf{File} & \textbf{Contents}\\
\midrule
\texttt{preamble.tex}
& Shared \LaTeX\ preamble: packages, colour definitions,
  tcolorbox styles, programme macros.
  Input by all artifact files via \texttt{\textbackslash input\{../preamble\}}.\\
\texttt{L01\_manifest.tex}
& This file. Complete inventory of the L01 artifact suite.\\
\bottomrule
\end{tabularx}

\bigskip
%=============================================================
\section*{Artifact Count and Coverage}
%=============================================================

\begin{center}
\begin{tabular}{llcc}
\toprule
\textbf{Code} & \textbf{Full name} & \textbf{Tiers} & \textbf{Delivered}\\
\midrule
SCQU & Simple Concept Questions UG      & 1--3 & \checkmark\\
CCQU & Complex Concept Questions UG     & 2--3 & \checkmark\\
SPSU & Simple Problem Set UG            & 1--3 & \checkmark\\
CPSU & Complex Problem Set UG           & 3    & \checkmark\\
SPJU & Simple Projects UG               & 2--3 & \checkmark\\
CPJU & Complex Projects UG              & 3    & \checkmark\\
\midrule
SCQG & Simple Concept Questions Grad    & 4--5 & \checkmark\\
CCQG & Complex Concept Questions Grad   & 4--5 & \checkmark\\
SPSG & Simple Problem Set Grad          & 4--5 & \checkmark\\
CPSG & Complex Problem Set Grad         & 4--5 & \checkmark\\
SPJG & Simple Projects Grad             & 4--5 & \checkmark\\
CPJG & Complex Projects Grad            & 5    & \checkmark\\
\midrule
WRQ  & Written Research Questions       & 4--5 & \checkmark\\
ORQ  & Oral Research Questions          & 5    & \checkmark\\
\midrule
\textbf{Total} & & & \textbf{14 / 14}\\
\bottomrule
\end{tabular}
\end{center}

\bigskip
%=============================================================
\section*{Compilation Instructions}
%=============================================================

\begin{tcolorbox}[goldbox]
\textbf{To compile any artifact file:}
\begin{verbatim}
  cd L01-artifacts/
  pdflatex ug/L01_SCQU.tex
  pdflatex ug/L01_CCQU.tex
  pdflatex ug/L01_SPSU.tex
  pdflatex ug/L01_CPSU.tex
  pdflatex ug/L01_SPJU.tex
  pdflatex ug/L01_CPJU.tex
  pdflatex grad/L01_SCQG.tex
  pdflatex grad/L01_CCQG.tex
  pdflatex grad/L01_SPSG.tex
  pdflatex grad/L01_CPSG.tex
  pdflatex grad/L01_SPJG_CPJG.tex
  pdflatex research/L01_WRQ_ORQ.tex
  pdflatex bibliography/L01_bibliography_additions.tex
  pdflatex lecture/L01_lecture_body.tex
  pdflatex L01_manifest.tex
\end{verbatim}
All files \texttt{\textbackslash input\{../preamble\}} so must be compiled from
their subdirectory, or the \texttt{\textbackslash input} path adjusted.
Run \texttt{pdflatex} twice for correct cross-references.
\end{tcolorbox}

\bigskip
%=============================================================
\section*{Relationship to Slide Deck}
%=============================================================

\begin{tcolorbox}[notebox]
\textbf{Important.}
All 14 artifact files in this suite are \textbf{separate from the slide deck}.
The slide deck (\texttt{slides/L01\_expanded\_Pure\_Mixed\_States\_and\_the\_Density\_Matrix.pptx})
covers the lecture content in 28 slides with diagrams for each learning objective,
formula, concept question, and tier. The artifacts are standalone documents
distributed to students separately according to their tier track.

The slide deck was generated by \texttt{src/builder-L01.js} and contains no
problem sets, reading lists, or assessment rubrics --- only lecture content diagrams.
\end{tcolorbox}

\end{document}
